% Options for packages loaded elsewhere
\PassOptionsToPackage{unicode}{hyperref}
\PassOptionsToPackage{hyphens}{url}
\PassOptionsToPackage{dvipsnames,svgnames,x11names}{xcolor}
\documentclass[
  11pt,
]{article}
\usepackage{xcolor}
\usepackage[margin=2cm]{geometry}
\usepackage{amsmath,amssymb}
\setcounter{secnumdepth}{-\maxdimen} % remove section numbering
\usepackage{iftex}
\ifPDFTeX
  \usepackage[T1]{fontenc}
  \usepackage[utf8]{inputenc}
  \usepackage{textcomp} % provide euro and other symbols
\else % if luatex or xetex
  \usepackage{unicode-math} % this also loads fontspec
  \defaultfontfeatures{Scale=MatchLowercase}
  \defaultfontfeatures[\rmfamily]{Ligatures=TeX,Scale=1}
\fi
\usepackage{lmodern}
\ifPDFTeX\else
  % xetex/luatex font selection
    \setmonofont[Scale=0.8,BoldFont=DejaVuSansMono-Bold.ttf]{DejaVuSansMono.ttf}
\fi
% Use upquote if available, for straight quotes in verbatim environments
\IfFileExists{upquote.sty}{\usepackage{upquote}}{}
\IfFileExists{microtype.sty}{% use microtype if available
  \usepackage[]{microtype}
  \UseMicrotypeSet[protrusion]{basicmath} % disable protrusion for tt fonts
}{}
\makeatletter
\@ifundefined{KOMAClassName}{% if non-KOMA class
  \IfFileExists{parskip.sty}{%
    \usepackage{parskip}
  }{% else
    \setlength{\parindent}{0pt}
    \setlength{\parskip}{6pt plus 2pt minus 1pt}}
}{% if KOMA class
  \KOMAoptions{parskip=half}}
\makeatother
\usepackage{color}
\usepackage{fancyvrb}
\newcommand{\VerbBar}{|}
\newcommand{\VERB}{\Verb[commandchars=\\\{\}]}
\DefineVerbatimEnvironment{Highlighting}{Verbatim}{commandchars=\\\{\}}
% Add ',fontsize=\small' for more characters per line
\newenvironment{Shaded}{}{}
\newcommand{\AlertTok}[1]{\textcolor[rgb]{1.00,0.00,0.00}{\textbf{#1}}}
\newcommand{\AnnotationTok}[1]{\textcolor[rgb]{0.38,0.63,0.69}{\textbf{\textit{#1}}}}
\newcommand{\AttributeTok}[1]{\textcolor[rgb]{0.49,0.56,0.16}{#1}}
\newcommand{\BaseNTok}[1]{\textcolor[rgb]{0.25,0.63,0.44}{#1}}
\newcommand{\BuiltInTok}[1]{\textcolor[rgb]{0.00,0.50,0.00}{#1}}
\newcommand{\CharTok}[1]{\textcolor[rgb]{0.25,0.44,0.63}{#1}}
\newcommand{\CommentTok}[1]{\textcolor[rgb]{0.38,0.63,0.69}{\textit{#1}}}
\newcommand{\CommentVarTok}[1]{\textcolor[rgb]{0.38,0.63,0.69}{\textbf{\textit{#1}}}}
\newcommand{\ConstantTok}[1]{\textcolor[rgb]{0.53,0.00,0.00}{#1}}
\newcommand{\ControlFlowTok}[1]{\textcolor[rgb]{0.00,0.44,0.13}{\textbf{#1}}}
\newcommand{\DataTypeTok}[1]{\textcolor[rgb]{0.56,0.13,0.00}{#1}}
\newcommand{\DecValTok}[1]{\textcolor[rgb]{0.25,0.63,0.44}{#1}}
\newcommand{\DocumentationTok}[1]{\textcolor[rgb]{0.73,0.13,0.13}{\textit{#1}}}
\newcommand{\ErrorTok}[1]{\textcolor[rgb]{1.00,0.00,0.00}{\textbf{#1}}}
\newcommand{\ExtensionTok}[1]{#1}
\newcommand{\FloatTok}[1]{\textcolor[rgb]{0.25,0.63,0.44}{#1}}
\newcommand{\FunctionTok}[1]{\textcolor[rgb]{0.02,0.16,0.49}{#1}}
\newcommand{\ImportTok}[1]{\textcolor[rgb]{0.00,0.50,0.00}{\textbf{#1}}}
\newcommand{\InformationTok}[1]{\textcolor[rgb]{0.38,0.63,0.69}{\textbf{\textit{#1}}}}
\newcommand{\KeywordTok}[1]{\textcolor[rgb]{0.00,0.44,0.13}{\textbf{#1}}}
\newcommand{\NormalTok}[1]{#1}
\newcommand{\OperatorTok}[1]{\textcolor[rgb]{0.40,0.40,0.40}{#1}}
\newcommand{\OtherTok}[1]{\textcolor[rgb]{0.00,0.44,0.13}{#1}}
\newcommand{\PreprocessorTok}[1]{\textcolor[rgb]{0.74,0.48,0.00}{#1}}
\newcommand{\RegionMarkerTok}[1]{#1}
\newcommand{\SpecialCharTok}[1]{\textcolor[rgb]{0.25,0.44,0.63}{#1}}
\newcommand{\SpecialStringTok}[1]{\textcolor[rgb]{0.73,0.40,0.53}{#1}}
\newcommand{\StringTok}[1]{\textcolor[rgb]{0.25,0.44,0.63}{#1}}
\newcommand{\VariableTok}[1]{\textcolor[rgb]{0.10,0.09,0.49}{#1}}
\newcommand{\VerbatimStringTok}[1]{\textcolor[rgb]{0.25,0.44,0.63}{#1}}
\newcommand{\WarningTok}[1]{\textcolor[rgb]{0.38,0.63,0.69}{\textbf{\textit{#1}}}}
\usepackage{longtable,booktabs,array}
\usepackage{calc} % for calculating minipage widths
% Correct order of tables after \paragraph or \subparagraph
\usepackage{etoolbox}
\makeatletter
\patchcmd\longtable{\par}{\if@noskipsec\mbox{}\fi\par}{}{}
\makeatother
% Allow footnotes in longtable head/foot
\IfFileExists{footnotehyper.sty}{\usepackage{footnotehyper}}{\usepackage{footnote}}
\makesavenoteenv{longtable}
\ifLuaTeX
\usepackage[bidi=basic]{babel}
\else
\usepackage[bidi=default]{babel}
\fi
\babelprovide[main,import]{french}
% get rid of language-specific shorthands (see #6817):
\let\LanguageShortHands\languageshorthands
\def\languageshorthands#1{}
\setlength{\emergencystretch}{3em} % prevent overfull lines
\providecommand{\tightlist}{%
  \setlength{\itemsep}{0pt}\setlength{\parskip}{0pt}}
% Ajouté au préambule de chaque document par generer.sh
\usepackage{xcolor}
\usepackage{newunicodechar}
\newfontfamily\symbolfont{DejaVuSans.ttf}
% Caractères absents de Latin Modern (flèches, comparateurs, cases) : repli sur DejaVu Sans
\newunicodechar{→}{{\symbolfont →}}
\newunicodechar{←}{{\symbolfont ←}}
\newunicodechar{⇄}{{\symbolfont ⇄}}
\newunicodechar{↔}{{\symbolfont ↔}}
\newunicodechar{≥}{{\symbolfont ≥}}
\newunicodechar{≤}{{\symbolfont ≤}}
\newunicodechar{≠}{{\symbolfont ≠}}
\newunicodechar{≈}{{\symbolfont ≈}}
\newunicodechar{✓}{{\symbolfont ✓}}
\newunicodechar{✗}{{\symbolfont ✗}}
\newunicodechar{⚠}{{\symbolfont ⚠}}
\newunicodechar{✅}{{\symbolfont\color[rgb]{0.1,0.55,0.2}✔}}
\newunicodechar{‑}{\mbox{-}}
% Retour à la ligne dans les blocs de code et les URL longues
\usepackage{fvextra}
\DefineVerbatimEnvironment{Highlighting}{Verbatim}{breaklines,commandchars=\\\{\}}
\usepackage{xurl}
\setlength{\emergencystretch}{3em}
\AtBeginEnvironment{longtable}{\small}
% Figures flottantes (pas de demi-pages vides), jamais plus hautes que 85 % de la page
\usepackage{float}
\floatplacement{figure}{htbp}
\usepackage{adjustbox}
\providecommand*\pandocbounded[1]{#1}% absent des documents sans image
\renewcommand*\pandocbounded[1]{\adjustbox{max width=\linewidth,max totalheight=0.85\textheight}{#1}}
\usepackage[font=small,labelfont=bf]{caption}
% Contenu sur une nouvelle page après la table des matières
\let\tdmorigine\tableofcontents
\renewcommand{\tableofcontents}{\tdmorigine\newpage}
\usepackage{fancyhdr}
\pagestyle{fancy}
\fancyhf{}
\fancyhead[L]{\small SmartRecruit — Qualité logicielle ISO/IEC 25010}
\fancyhead[R]{\small\thepage}
\renewcommand{\headrulewidth}{0.4pt}
\usepackage{bookmark}
\IfFileExists{xurl.sty}{\usepackage{xurl}}{} % add URL line breaks if available
\urlstyle{same}
\hypersetup{
  pdftitle={Fiche d\textquotesingle étude --- Membre 2 : interface utilisateur, E2E, compatibilité, accessibilité},
  pdflang={fr},
  colorlinks=true,
  linkcolor={Maroon},
  filecolor={Maroon},
  citecolor={Blue},
  urlcolor={Blue},
  pdfcreator={LaTeX via pandoc}}

\title{Fiche d\textquotesingle étude --- Membre 2 : interface
utilisateur, E2E, compatibilité, accessibilité}
\author{}
\date{}

\begin{document}
\maketitle

\renewcommand*\contentsname{Table des matières}
{
\hypersetup{linkcolor=}
\setcounter{tocdepth}{3}
\tableofcontents
}
\begin{quote}
Pour : \emph{(à compléter)}. Complète le \href{SOUTENANCE.pdf}{guide de
soutenance} (questions-réponses, membre 2) et la
\href{REPARTITION_TESTS.pdf}{répartition des tests}.
\end{quote}

\section{1. Ta mission en soutenance}\label{1-ta-mission-en-soutenance}

Le lead lance Playwright en direct ; toi, tu dois pouvoir
\textbf{expliquer ce qui s\textquotesingle exécute}, le
\textbf{relancer} si l\textquotesingle enseignant le demande et
\textbf{répondre} sur :

\begin{itemize}
\tightlist
\item
  \textbf{6.1 Aptitude fonctionnelle} (côté parcours utilisateur, tests
  E2E) ;
\item
  \textbf{6.3 Compatibilité} (≥ 2 navigateurs) ;
\item
  \textbf{6.4 Facilité d\textquotesingle utilisation} (parcours validés,
  accessibilité de base) ;
\item
  les tests unitaires du frontend (Vitest).
\end{itemize}

\section{2. Ce que tu dois apprendre
(notions)}\label{2-ce-que-tu-dois-apprendre-notions}

\begin{longtable}[]{@{}
  >{\raggedright\arraybackslash}p{(\linewidth - 2\tabcolsep) * \real{0.3860}}
  >{\raggedright\arraybackslash}p{(\linewidth - 2\tabcolsep) * \real{0.6140}}@{}}
\toprule\noalign{}
\begin{minipage}[b]{\linewidth}\raggedright
Notion
\end{minipage} & \begin{minipage}[b]{\linewidth}\raggedright
Ce que tu dois savoir dire
\end{minipage} \\
\midrule\noalign{}
\endhead
\bottomrule\noalign{}
\endlastfoot
Test E2E (de bout en bout) & Un vrai navigateur pilote la vraie
application (frontend + gateway + services + bases), comme un
utilisateur. Lent mais le plus proche du réel. \\
Test de composant (Vitest + Testing Library) & Un composant React rendu
seul, l\textquotesingle API simulée ; rapide, vérifie
l\textquotesingle affichage et les réactions (filtres, erreurs). \\
Playwright & Outil E2E : un fichier de config, des \textbf{projets} (un
par navigateur), des \textbf{tests} dans \texttt{tests/}, un rapport
HTML. \\
Projet \texttt{setup} et dépendances & \texttt{auth.setup.ts} se
connecte une fois par rôle et enregistre la session ; les autres projets
en dépendent et la réutilisent. \\
Auto-attente & Playwright attend qu\textquotesingle un élément soit
visible/actionnable avant d\textquotesingle agir : pas de
\texttt{sleep}, moins de tests instables. \\
Localisateurs accessibles & \texttt{getByRole}, \texttt{getByLabel} : on
cible ce que voit l\textquotesingle utilisateur (rôle, libellé), pas une
classe CSS fragile. \\
Test instable (« flaky ») & Un test qui passe ou échoue sans changement
de code (attentes trop courtes, données partagées). \\
Moteurs de navigateur & Chromium (Chrome, Edge), WebKit (Safari), Gecko
(Firefox). Tester deux moteurs différents a plus de valeur que deux
navigateurs Chromium. \\
Émulation mobile & Playwright simule l\textquotesingle écran, le tactile
et l\textquotesingle agent utilisateur d\textquotesingle un Pixel 7 ; ce
n\textquotesingle est pas un vrai téléphone. \\
Accessibilité, WCAG 2.1 A/AA & Règles pour que
l\textquotesingle application soit utilisable par tous (contrastes,
libellés, clavier, structure). A = minimum, AA = niveau visé
habituellement. \\
axe-core & Moteur qui vérifie automatiquement une partie des règles WCAG
sur une page ; il ne remplace pas un test manuel. \\
Contrôle d\textquotesingle accès côté interface & \texttt{RequireRole}
cache les pages selon le rôle ; la vraie sécurité est côté serveur (un
appel API sans le bon rôle reçoit 403). \\
\end{longtable}

\section{3. Fichiers à lire, dans
l\textquotesingle ordre}\label{3-fichiers-uxe0-lire-dans-lordre}

\begin{longtable}[]{@{}
  >{\raggedright\arraybackslash}p{(\linewidth - 4\tabcolsep) * \real{0.0237}}
  >{\raggedright\arraybackslash}p{(\linewidth - 4\tabcolsep) * \real{0.5621}}
  >{\raggedright\arraybackslash}p{(\linewidth - 4\tabcolsep) * \real{0.4142}}@{}}
\toprule\noalign{}
\begin{minipage}[b]{\linewidth}\raggedright
\#
\end{minipage} & \begin{minipage}[b]{\linewidth}\raggedright
Fichier
\end{minipage} & \begin{minipage}[b]{\linewidth}\raggedright
Ce que tu y cherches
\end{minipage} \\
\midrule\noalign{}
\endhead
\bottomrule\noalign{}
\endlastfoot
1 &
\texttt{quality/}\allowbreak{}\texttt{e2e-}\allowbreak{}\texttt{ui/}\allowbreak{}\texttt{playwright.}\allowbreak{}\texttt{config.}\allowbreak{}\texttt{ts}
& Les projets : \texttt{setup}, \texttt{chromium}, \texttt{edge},
\texttt{webkit}, \texttt{mobile-chrome}, et \texttt{firefox} seulement
si \texttt{PW\_FIREFOX} est défini. \\
2 &
\texttt{quality/}\allowbreak{}\texttt{e2e-}\allowbreak{}\texttt{ui/}\allowbreak{}\texttt{tests/}\allowbreak{}\texttt{accounts.}\allowbreak{}\texttt{ts}
et \texttt{auth.setup.ts} & Les comptes de démo et la connexion
préalable des rôles. \\
3 &
\texttt{quality/}\allowbreak{}\texttt{e2e-}\allowbreak{}\texttt{ui/}\allowbreak{}\texttt{tests/}\allowbreak{}\texttt{journeys.}\allowbreak{}\texttt{spec.}\allowbreak{}\texttt{ts}
& Les \textbf{13 parcours} : savoir les citer et expliquer ce que
vérifie chacun. \\
4 &
\texttt{quality/}\allowbreak{}\texttt{e2e-}\allowbreak{}\texttt{ui/}\allowbreak{}\texttt{tests/}\allowbreak{}\texttt{a11y.}\allowbreak{}\texttt{spec.}\allowbreak{}\texttt{ts}
& Les \textbf{6 pages} auditées par axe-core et les règles
appliquées. \\
5 &
\texttt{frontend/}\allowbreak{}\texttt{src/}\allowbreak{}\texttt{App.}\allowbreak{}\texttt{tsx}
et
\texttt{frontend/}\allowbreak{}\texttt{src/}\allowbreak{}\texttt{auth/}\allowbreak{}\texttt{Require}\allowbreak{}\texttt{Role.}\allowbreak{}\texttt{tsx}
& Les routes et la protection par rôle (redirection, page
\texttt{/forbidden}). \\
6 &
\texttt{frontend/}\allowbreak{}\texttt{src/}\allowbreak{}\texttt{api/}\allowbreak{}\texttt{client.}\allowbreak{}\texttt{ts},
\texttt{frontend/}\allowbreak{}\texttt{src/}\allowbreak{}\texttt{api/}\allowbreak{}\texttt{mock/}\allowbreak{}
& Comment l\textquotesingle interface appelle l\textquotesingle API, et
la simulation utilisée par les tests Vitest. \\
7 & Deux tests Vitest, par ex.
\texttt{frontend/}\allowbreak{}\texttt{src/}\allowbreak{}\texttt{auth/}\allowbreak{}\texttt{Require}\allowbreak{}\texttt{Role.}\allowbreak{}\texttt{test.}\allowbreak{}\texttt{tsx}
et un test de page dans \texttt{frontend/src/pages/} & Structure
d\textquotesingle un test de composant : rendu, action, assertion. \\
8 & \href{MATRICE_TRACABILITE.pdf}{MATRICE\_TRACABILITE.md}, lignes
EF‑01 à EF‑06, EQ‑01, EQ‑03, EQ‑04 & Où tes tests apparaissent comme
preuve. \\
\end{longtable}

\section{4. Commandes à savoir lancer (et
lire)}\label{4-commandes-uxe0-savoir-lancer-et-lire}

Application démarrée au préalable
(\texttt{docker\ compose\ up\ -}\allowbreak{}\texttt{d\ -}\allowbreak{}\texttt{-}\allowbreak{}\texttt{build},
données de démo chargées).

\begin{Shaded}
\begin{Highlighting}[]
\FunctionTok{bash}\NormalTok{ quality/run{-}all.sh ui}
\end{Highlighting}
\end{Shaded}

Toute la suite (4 navigateurs + accessibilité). Lire :
\texttt{62\ passed} (et des tests \texttt{skipped}, voir section 7).

\begin{Shaded}
\begin{Highlighting}[]
\BuiltInTok{cd}\NormalTok{ quality/e2e{-}ui }\KeywordTok{\&\&} \ExtensionTok{npx}\NormalTok{ playwright test }\AttributeTok{{-}{-}project}\OperatorTok{=}\NormalTok{chromium }\AttributeTok{{-}{-}headed}
\end{Highlighting}
\end{Shaded}

Un seul navigateur, fenêtre visible : la version de la démo.

\begin{Shaded}
\begin{Highlighting}[]
\BuiltInTok{cd}\NormalTok{ quality/e2e{-}ui }\KeywordTok{\&\&} \ExtensionTok{npx}\NormalTok{ playwright show{-}report ../reports/e2e{-}ui/html}
\end{Highlighting}
\end{Shaded}

Ouvre le rapport HTML : savoir montrer un test, ses étapes et sa durée.

\begin{Shaded}
\begin{Highlighting}[]
\BuiltInTok{cd}\NormalTok{ frontend }\KeywordTok{\&\&} \ExtensionTok{npm}\NormalTok{ test}
\end{Highlighting}
\end{Shaded}

Tests Vitest du frontend (403 tests).

\section{5. Chiffres à connaître par
cœur}\label{5-chiffres-uxe0-connauxeetre-par-cux153ur}

\begin{longtable}[]{@{}
  >{\raggedright\arraybackslash}p{(\linewidth - 2\tabcolsep) * \real{0.3171}}
  >{\raggedright\arraybackslash}p{(\linewidth - 2\tabcolsep) * \real{0.6829}}@{}}
\toprule\noalign{}
\begin{minipage}[b]{\linewidth}\raggedright
Sujet
\end{minipage} & \begin{minipage}[b]{\linewidth}\raggedright
Valeur
\end{minipage} \\
\midrule\noalign{}
\endhead
\bottomrule\noalign{}
\endlastfoot
Parcours E2E & \textbf{13}, sur \textbf{4} navigateurs/appareils
(Chromium, Edge, WebKit, Chrome mobile Pixel 7) \\
Total Playwright & \textbf{62 réussis} = 13 × 4 + 6 audits
d\textquotesingle accessibilité (Chromium) + 4 connexions de
préparation \\
Accessibilité & axe-core 4.10, WCAG 2.0/2.1 A et AA, \textbf{6 pages},
\textbf{0 violation} \\
Vitest & \textbf{403} tests, 32 fichiers, couverture frontend
\textbf{92,1 \%} des lignes \\
Firefox & \textbf{non exécutable} sur le poste (binaire bloqué par
Windows, \texttt{spawn\ UNKNOWN}) : documenté, pas simulé \\
\end{longtable}

\section{6. Preuves à savoir ouvrir}\label{6-preuves-uxe0-savoir-ouvrir}

\begin{itemize}
\tightlist
\item
  \texttt{quality/}\allowbreak{}\texttt{reports/}\allowbreak{}\texttt{e2e-}\allowbreak{}\texttt{ui/}\allowbreak{}\texttt{html/}\allowbreak{}\texttt{index.}\allowbreak{}\texttt{html}
  (rapport Playwright) et
  \texttt{quality/}\allowbreak{}\texttt{reports/}\allowbreak{}\texttt{e2e-}\allowbreak{}\texttt{ui/}\allowbreak{}\texttt{junit.}\allowbreak{}\texttt{xml}.
\item
  \texttt{quality/}\allowbreak{}\texttt{reports/}\allowbreak{}\texttt{a11y/}\allowbreak{}\texttt{*.}\allowbreak{}\texttt{json}
  : un fichier par page, tableau de violations vide.
\item
  \texttt{quality/}\allowbreak{}\texttt{reports/}\allowbreak{}\texttt{unit/}\allowbreak{}\texttt{frontend-}\allowbreak{}\texttt{junit.}\allowbreak{}\texttt{xml}
  et
  \texttt{quality/}\allowbreak{}\texttt{reports/}\allowbreak{}\texttt{coverage/}\allowbreak{}\texttt{frontend/}\allowbreak{}.
\end{itemize}

\section{7. Questions pièges (au-delà du guide de
soutenance)}\label{7-questions-piuxe8ges-au-deluxe0-du-guide-de-soutenance}

\begin{longtable}[]{@{}
  >{\raggedright\arraybackslash}p{(\linewidth - 2\tabcolsep) * \real{0.4667}}
  >{\raggedright\arraybackslash}p{(\linewidth - 2\tabcolsep) * \real{0.5333}}@{}}
\toprule\noalign{}
\begin{minipage}[b]{\linewidth}\raggedright
Question
\end{minipage} & \begin{minipage}[b]{\linewidth}\raggedright
Réponse
\end{minipage} \\
\midrule\noalign{}
\endhead
\bottomrule\noalign{}
\endlastfoot
Pourquoi des tests « skipped » dans le rapport ? & Les audits
d\textquotesingle accessibilité ne tournent que sur Chromium : axe
analyse le DOM, qui est le même quel que soit le navigateur. Les répéter
n\textquotesingle apporterait rien. \\
Edge et Chrome, c\textquotesingle est le même moteur : est-ce vraiment 2
navigateurs ? & Oui, deux navigateurs, mais le même moteur.
C\textquotesingle est pour ça qu\textquotesingle on teste aussi
\textbf{WebKit} (moteur de Safari), un moteur réellement différent. \\
Le mobile est émulé : est-ce un vrai test de compatibilité ? &
C\textquotesingle est une émulation (taille d\textquotesingle écran,
tactile, agent utilisateur) : elle valide la mise en page et les
parcours, pas les particularités d\textquotesingle un vrai appareil.
Limite assumée. \\
Que se passe-t-il si un test E2E échoue ? & Playwright garde une capture
d\textquotesingle écran (toujours) et une trace complète (en cas
d\textquotesingle échec) dans
\texttt{quality/}\allowbreak{}\texttt{reports/}\allowbreak{}\texttt{e2e-}\allowbreak{}\texttt{ui/}\allowbreak{}\texttt{artifacts}
: la trace se rejoue pas à pas (DOM, réseau, console) pour trouver la
cause. \\
Pourquoi ne pas tout tester en E2E ? & Trop lent et fragile ; les
détails (filtres, messages d\textquotesingle erreur, cas limites) sont
testés par Vitest, l\textquotesingle E2E vérifie que les parcours
complets marchent ensemble. \\
Comment l\textquotesingle E2E respecte-t-il l\textquotesingle anti-force
brute (RG‑06) ? & La connexion est faite une seule fois par rôle
(\texttt{auth.setup.ts}) puis réutilisée : 4 connexions par exécution,
sous la limite de 10 par 5 minutes. \\
La facilité d\textquotesingle utilisation, c\textquotesingle est
seulement axe ? & Non : parcours utilisateur validés de bout en bout +
choix de conception (libellés visibles, lien « Aller au contenu »,
navigation clavier, messages d\textquotesingle erreur clairs, mise en
page mobile). \\
Masquer un bouton selon le rôle, c\textquotesingle est de la sécurité ?
& Non, c\textquotesingle est du confort. Le serveur refuse
l\textquotesingle action (403) même si l\textquotesingle appel est forgé
; testé côté backend et par le parcours \texttt{/forbidden}. \\
\end{longtable}

\section{8. Plan
d\textquotesingle apprentissage}\label{8-plan-dapprentissage}

\begin{longtable}[]{@{}
  >{\raggedright\arraybackslash}p{(\linewidth - 4\tabcolsep) * \real{0.1316}}
  >{\raggedright\arraybackslash}p{(\linewidth - 4\tabcolsep) * \real{0.1316}}
  >{\raggedright\arraybackslash}p{(\linewidth - 4\tabcolsep) * \real{0.7368}}@{}}
\toprule\noalign{}
\begin{minipage}[b]{\linewidth}\raggedright
Étape
\end{minipage} & \begin{minipage}[b]{\linewidth}\raggedright
Durée
\end{minipage} & \begin{minipage}[b]{\linewidth}\raggedright
Objectif
\end{minipage} \\
\midrule\noalign{}
\endhead
\bottomrule\noalign{}
\endlastfoot
1 & 1 h & Utiliser l\textquotesingle application avec les 3 comptes de
démo : refaire à la main les 13 parcours. \\
2 & 2 h & Lire
\texttt{playwright.}\allowbreak{}\texttt{config.}\allowbreak{}\texttt{ts},
\texttt{auth.setup.ts}, \texttt{journeys.spec.ts} ; relier chaque test
au parcours fait à la main. \\
3 & 1 h & Lancer \texttt{ui}, ouvrir le rapport HTML et un fichier
\texttt{a11y/*.json}. \\
4 & 1 h & Lire \texttt{RequireRole.tsx} et deux tests Vitest ; lancer
\texttt{npm\ test}. \\
5 & 1 h & Réciter les réponses du guide de soutenance (membre 2) et de
la section 7. \\
\end{longtable}

\section{9. Auto-évaluation}\label{9-auto-uxe9valuation}

\begin{itemize}
\tightlist
\item[$\square$]
  Je cite les 13 parcours E2E et les 4 navigateurs.
\item[$\square$]
  J\textquotesingle explique le calcul 62 = 13 × 4 + 6 + 4, et les tests
  « skipped ».
\item[$\square$]
  J\textquotesingle explique pourquoi Firefox n\textquotesingle a pas de
  résultat, et la procédure pour l\textquotesingle exécuter ailleurs.
\item[$\square$]
  Je lance la suite et j\textquotesingle ouvre le rapport HTML sans
  aide.
\item[$\square$]
  J\textquotesingle explique la différence entre un test Vitest et un
  test Playwright, avec un exemple de chaque.
\item[$\square$]
  Je dis ce que vérifie axe-core et ce qu\textquotesingle il ne vérifie
  pas.
\end{itemize}

\end{document}
